\subsection{Introduction}
When no limit is placed on the number of warping fields \(\bm{\alpha}\) employed by a beam theory, it is possible to represent a continuum to any desired accuracy. 
However this is impractical, and it is preferable to forgo auxiliary warping fields altogether whenever possible. 
In the case of linear isotropy this is achieved by employing cross sectional constants like the shear correction factor. However, even in the simple case of linear isotropy, the subject of shear warping in beam theory remains a actively debated subject with unresolved ambiguities and strong disagreements \citep{dong2010much,dong2013principal,serpieri2014equivalence}. 

The difficulty may be attributed to two issues:
\begin{enumerate}
\item Indeterminacy in Saint~Venant's continuum solutions, and
\item Indeterminacy in the reduced form
\RepeatEquation{eq:trace-lift-linear}
% \begin{equation*}
% % \begin{SaveEquation}{eq:trace-lift-linear}
% \hat{\bm u}(\reflenx, \bm{r})=\bm{u}(x)
% +\bm\theta(x)\times\bm{r}+\bm\Phi(\bm{r})\,\bm{\alpha}(x)
% % \label{eq:trace-lift-linear}
% % \end{SaveEquation}
% \end{equation*}
\end{enumerate}
%
The ambiguous nature of Saint~Venant's solutions is well known in the continuum mechanics literature (\cref{sec:sv}). 
However, we are not aware of any works that formalize this in the context of dimensionally reduced 1D beam theories, or systematically investigate the extent to which beam theories may exactly represent 3D solutions in the Saint Venant class. 
Rather, it is typically assumed at the outset that the 1D model will only approximate a continuum solution. 
The subject is raised by \cite{simo1982bifurcation}, where a brief investigation is made in the context of planar buckling. 
However, because this investigation was limited to the plane, several salient three dimensional considerations do not arise such as work-conjugacy of resultants and handling of asymmetric sections. 

% The representation \eqref{eq:trace-lift-linear} is non-unique, which
% poses a fundamental obstacle for associating these modes to classical deflection \(\bm{u}(\reflenx)\), rotation \(\bm{\theta}(\reflenx)\), and warping \(\bm{\alpha}(\reflenx)\) fields. 


\begin{itemize}
    \item \cite{timoshenko1987theory} investigated a the flexural solution identified by the root boundary condition:
    \[
    \hat{\bm u}'(x_0, \mathbf{0}) = \mathbf{0}
    \]
    A trace defined to realize boundary conditions in this form exactly reproduces the response of classical Euler-Bernoulli beam theory, while remaining exact in the Saint-Venant class.

    \item \cite{timoshenko1921correction} imposed the root condition by requiring the directional derivative of the axial displacement component along a chosen cross-section direction \(\mathbf{t}\) vanishes at that boundary point. 
    In 3D this might be stated as:
    $$
    \left.\mathbf{t} \cdot \nabla_r(\hat{\bm u} \cdot \FrameBasis)\right|_{x=L, \bm{r}=r_0}=0,
    $$
    This produces a shear correction of \(\kappa = 2/3\).
    % The full trace is derived in \cref{sec:trace-timo}.
    \item In the landmark work by \cite{cowper1966shear}, a systematic investigation was conducted to obtain shear-correction factors corresponding to a trace which we have generalized to 3D with the form:
    \begin{equation}
    \bar{\bm{u}}(x) \coloneq \SectionAverage{\hat{\bm{u}}} - \bar{\bm{\theta}} \times \CentroidLocation, 
    \qquad 
    \bar{\bm{\theta}}(x) \coloneq \RotationAverage{\hat{\bm{u}}},
    \label{eq:def-section-avg-trace}
    \end{equation}
    Here $\SectionAverage{\cdot}$ and $\RotationAverage{\cdot}$ are integral functionals defined by \eqref{eq:def-section-averages} which average the 3D field \(\hat{\bm u}\) at a cross section and furnish an $L^2$-optimal fitting of a section-rigid field to the 3D displacement \(\hat{\bm u}_{\rm P}\). 
    The free parameters are determined by section geometry through rotation and translation averages of the warping functions, which is carried out in \cref{sec:trace-cowper}.
    The approach by \cite{cowper1966shear} laid the foundation for the plane beam theories developed by \cite{simo1982bifurcation,hjelmstad1983seismic,hjelmstad1987warping}. 

    \item The shear-correction factors of \cite{krenk1985pretwist}, \cite{renton1991generalized} and \cite{romano1992shear} are realized by a trace that defines beam strains to be 
    work-conjugate to the classical resultants:
    \begin{equation}
    \delta\bm{\gamma} \cdot \bm{n} + \delta\bm{\kappa} \cdot \bm{m} 
    = \int_\Section \delta\bm{\varepsilon} \cdot \bm{\sigma} \, \mathrm{d}A
    \label{eq:def-trace-conjugate}
    \end{equation}
    for all variations within the Saint~Venant class.
      % \begin{equation*}
      %   \int \ShearStress_{(b)}\cdot (\ShearStress_{(b)}/G)\,\mathrm{d}A
      %   = \ShearForce\cdot \bm\gamma^{\rm sh} + T_{\mathrm C^{\rm sh}}\varkappa_{(b)}
      % \end{equation*}
    This uniquely determines all 18 
    parameters and guarantees a symmetric section stiffness matrix. The full trace is derived in \cref{sec:trace-energy}. 
    \item
    \cite[Eq. (7.4.52)]{slaughter2013linearized}, and \cite{reddy2002energy}:
    \[
    \tfrac12 \bm{\epsilon}\nabla\times\hat{\bm u}= \mathbf{0}
    \]
    This produces a \emph{stiffening} shear correction of \(\kappa = 4/3\), which is peculiarly greater than unity. 
    % The full trace is derived in \cref{sec:trace-slaughter}.
\end{itemize}


\subsection{Outline}


An outline:
\begin{itemize}
    \item Introduce the \emph{trace} abstraction, which formalizes the ambiguity in representing the Saint-Venant class through classical fields \(\bm{u},\bm{\theta}\). 
    \item Investigate the special class of traces that exhibit Type 3 locality... they can be exactly represented with \emph{only} 6-DOF, and are compatible with the interpolation assumptions of traditional elements.
    \item Elaborate the effect of the trace on the boundary conditions that the beam model is capable of representing. 
\end{itemize}

